\documentclass[11pt,a4paper]{article}
\usepackage{curso}
\begin{document}
\portada{Evaluación --- Unidad}{5}{Quiz: hardware probabilístico y termodinámico}
\noindent Preguntas guía: ¿qué hace un chip de pbits mejor que uno digital? ¿Qué lo limita? ¿Cómo se cuenta honestamente su energía?
\section*{Preguntas}
\begin{enumerate}
\item Un pbit con entrada $u$ produce $m=+1$ con probabilidad: (a) $u$ (b) $\sigma(2\beta u)$ (c) $\tanh(u)$ (d) $1/2$ siempre.
\item La cota del Z1 por actualización ($\le74$ fJ) está a unas \rule{1.2cm}{0.4pt} veces del límite de Landauer a 300 K: (a) $10^2$ (b) $10^4$ (c) $2{,}6\times10^7$ (d) $10^{12}$.
\item Estacionaria de $dx=-(Ax-b)\,dt+\sqrt{2/\beta}\,dW$: (a) $\N(b,A)$ (b) $\N(A^{-1}b,\beta^{-1}A^{-1})$ (c) uniforme (d) $\N(0,I)$.
\item En un circuito browniano de tokens, sesgar el avance ($k_+>k_-$): (a) no cuesta nada (b) reduce el tiempo y cuesta disipación (c) aumenta la varianza relativa (d) requiere reloj.
\item Para desajustes pequeños de los campos, la KL crece como: (a) $\sigma$ (b) $\sigma^2$ (c) $e^{\sigma}$ (d) no depende de $\sigma$.
\item El redondeo estocástico es: (a) determinista (b) sesgado hacia cero (c) insesgado (d) solo para enteros.
\item El grafo padre de un chip es: (a) completo (b) fijo y disperso (c) elegido por el programa (d) aleatorio en cada barrido.
\item En el ejemplo de la lección 5.8 con pegamento fp32, el término dominante es: (a) tejido (b) E/S (c) pegamento digital (d) Landauer.
\item (Problema) ¿Cuántas copias mínimas necesita una variable de grado 50 en un grafo de grado 16?
\item (Problema) Probabilidad de ruptura de una cadena de dos copias con $\beta J_c=1{,}5$.
\item (Problema) Energía del tejido con $N=10^5$ pbits, $K=500$ barridos y 74 fJ por actualización.
\item (Problema) Varianza del redondeo estocástico de $x=q_k+0{,}3\Delta$.
\end{enumerate}
\newpage
\section*{Clave}
1 (b) · 2 (c) · 3 (b) · 4 (b) · 5 (b) · 6 (c) · 7 (b) · 8 (c).
\begin{enumerate}\setcounter{enumi}{8}
\item $\lceil(50-2)/(16-2)\rceil=4$.
\item $1/(1+e^{3})=0{,}047$.
\item $10^5\times500\times74\,\text{fJ}=3{,}7$ \textmu J \evE.
\item $\Delta^2\cdot0{,}3\cdot0{,}7=0{,}21\,\Delta^2$.
\end{enumerate}
Aprobación 8/12; fallos en 3, 5 o 8 $\Rightarrow$ repasar 5.4, 5.6, 5.8.
\end{document}
